
\documentclass[12pt, a4paper]{article}

\usepackage{listings}
\usepackage{xcolor}
\usepackage{fontspec}

\definecolor{codegreen}{rgb}{0,0.6,0}
\definecolor{codegray}{rgb}{0.5,0.5,0.5}
\definecolor{codepurple}{rgb}{0.58,0,0.82}
\definecolor{backcolour}{rgb}{0.95,0.95,0.92}

\lstdefinestyle{mystyle}{
  backgroundcolor=\color{backcolour},
  commentstyle=\color{codegreen},
  keywordstyle=\color{magenta},
  numberstyle=\tiny\color{codegray},
  stringstyle=\color{codepurple},
  basicstyle=\ttfamily\footnotesize,
  breaklines=true,
  keepspaces=true,
  numbers=left,
  numbersep=5pt,
  showstringspaces=false,
  tabsize=2
}

\title{Homework: Infix to Postfix}
\author{LEE JYUN WEI \hspace{0.1pt} D1450784}
\date{}

\lstset{style=mystyle}

\begin{document}
\maketitle

\section{Problem}

Read multiple infix expressions from
\texttt{1117.txt}.

Convert each infix expression into a postfix
expression using a stack, and calculate the
final result.

The program should also handle parentheses
and operator precedence.

\section{Code Snippet}

\lstinputlisting[language=C++]{prefix-postfix.cpp}

\section{Result}

Input:
\begin{verbatim}
4*2+10-9/3
\end{verbatim}

Output:
\begin{verbatim}
42*10+93/-
result : 15
\end{verbatim}

\section{Editorial}

The program uses two stacks to solve the problem.

First, a stack is used to convert the infix
expression into postfix notation.

Numbers are added directly to the postfix
expression. Operators are pushed onto the
stack according to their precedence.

Parentheses are handled separately to ensure
the correct order of operations.

After conversion, another stack is used to
calculate the postfix expression.

When an operator is found, the program pops
two numbers from the stack, performs the
calculation, and pushes the result back.

Finally, the remaining number is the answer.

\end{document}
